\chapter[APÊNDICE G -- RESUMO DAS RESPOSTAS AOS FORMULÁRIOS]{{\centering APÊNDICE G \textendash{} RESUMO DAS RESPOSTAS AOS FORMULÁRIOS}}  \label{ap:respostas}

Este apêndice reúne o resumo agregado das respostas aos dois formulários de avaliação, descritos na Seção~\ref{sec:protocolo-avaliacao} e reproduzidos nos Apêndices~\ref{ap:formulario-especialistas} e~\ref{ap:formulario-estudantes}. O resumo foi gerado por um \textit{script} a partir das planilhas exportadas dos formulários e contém apenas contagens, médias, medianas e valores mínimos e máximos. As planilhas, que trazem o horário de cada envio e os textos livres, não são publicadas. Os blocos de avaliação dos recursos, de uso em sala, de potencial pedagógico, de utilidade e intenção de uso e de riscos estão nas Tabelas~\ref{tab:resultado-recursos}, \ref{tab:resultado-ava}, \ref{tab:resultado-pedagogico} e~\ref{tab:resultado-riscos}, e as afirmações do formulário dos estudantes, na Tabela~\ref{tab:resultado-estudantes}; as tabelas deste apêndice completam esses dados.

Nas tabelas, ``n'' é o número de respostas dentro da escala. As respostas fora da escala (``não sei'', ``não usei'' ou ``não se aplica'') não entram na média e aparecem em coluna própria.

\section{Formulário dos Especialistas}

O formulário recebeu oito respostas entre 1º e 9 de novembro de 2026. A Tabela~\ref{tab:ap-perfil} traz as respostas às perguntas de perfil e de uso da plataforma. Nas perguntas que admitem mais de uma marcação, a soma passa de oito.

\begin{table}[htbp]
\centering
\caption{Perfil dos respondentes do formulário dos especialistas (oito respondentes).}
\label{tab:ap-perfil}
\footnotesize
\renewcommand{\arraystretch}{1.2}
\begin{tabular}{|p{10.5cm}|c|}
\hline
\textbf{Resposta} & \textbf{Respondentes} \\\hline
\multicolumn{2}{|l|}{\textbf{Maior titulação}} \\\hline
Ensino médio & 3 \\\hline
Doutorado ou superior & 3 \\\hline
Graduação & 2 \\\hline
\multicolumn{2}{|l|}{\textbf{Atuação (mais de uma marcação)}} \\\hline
Profissional de desenvolvimento de software & 5 \\\hline
Professor(a) do ensino superior & 3 \\\hline
Professor(a) do ensino técnico & 2 \\\hline
\multicolumn{2}{|l|}{\textbf{Ferramentas de ensino já conhecidas (mais de uma marcação)}} \\\hline
Portugol Studio ou Portugol Webstudio & 5 \\\hline
Scratch, Blockly ou similar & 4 \\\hline
Beecrowd, URI ou outro juiz online & 3 \\\hline
VisuAlg & 2 \\\hline
Thonny, Python Tutor ou outro visualizador de execução & 2 \\\hline
Nenhuma das anteriores & 1 \\\hline
\multicolumn{2}{|l|}{\textbf{Leciona ou já lecionou disciplina que envolva programação}} \\\hline
Não & 4 \\\hline
Sim, leciono ou já lecionei & 4 \\\hline
\multicolumn{2}{|l|}{\textbf{Tempo de trabalho com programação}} \\\hline
De 2 a 5 anos & 3 \\\hline
Menos de 2 anos & 2 \\\hline
De 6 a 10 anos & 2 \\\hline
Mais de 10 anos & 1 \\\hline
\multicolumn{2}{|l|}{\textbf{Tempo de uso da plataforma}} \\\hline
Menos de 15 minutos & 4 \\\hline
Mais de 1 hora & 2 \\\hline
De 30 a 60 minutos & 1 \\\hline
De 15 a 30 minutos & 1 \\\hline
\end{tabular}
\legend{Fonte: dados da pesquisa (2026).}
\end{table}

O que os respondentes declararam ter feito na plataforma e o que personalizaram está na Tabela~\ref{tab:resultado-atividades}. A Tabela~\ref{tab:ap-sus} traz os dez itens do SUS e o resumo do escore. Nos itens pares, de redação negativa, a concordância baixa é o resultado favorável.

\begin{table}[htbp]
\centering
\caption{Itens do SUS (escala de 1 a 5) e escore final; desvio padrão do escore: 16,5.}
\label{tab:ap-sus}
\footnotesize
\renewcommand{\arraystretch}{1.2}
\begin{tabular}{|p{8.4cm}|c|c|c|c|}
\hline
\textbf{Item} & \textbf{n} & \textbf{Média} & \textbf{Mediana} & \textbf{Mín.--máx.} \\\hline
1. Acho que gostaria de usar esta plataforma com frequência. & 8 & 3,50 & 3 & 2--5 \\\hline
2. Considerei a plataforma mais complexa do que precisava ser. & 8 & 1,88 & 2 & 1--3 \\\hline
3. Achei a plataforma fácil de usar. & 8 & 3,75 & 4 & 2--5 \\\hline
4. Acho que precisaria da ajuda de uma pessoa com conhecimento técnico para usar esta plataforma. & 8 & 1,88 & 2 & 1--3 \\\hline
5. Considerei que as várias funções da plataforma estavam bem integradas. & 8 & 4,25 & 4,5 & 3--5 \\\hline
6. Achei que havia muita inconsistência nesta plataforma. & 8 & 1,62 & 1,5 & 1--3 \\\hline
7. Imagino que a maioria das pessoas aprenderia a usar esta plataforma rapidamente. & 8 & 3,62 & 3,5 & 3--5 \\\hline
8. Achei a plataforma muito complicada de usar. & 8 & 1,75 & 2 & 1--3 \\\hline
9. Senti-me confiante ao usar a plataforma. & 8 & 3,88 & 4 & 3--5 \\\hline
10. Precisei aprender muitas coisas antes de conseguir usar esta plataforma. & 8 & 1,62 & 1,5 & 1--3 \\\hline
\textbf{Escore SUS (0 a 100)} & 8 & 75,6 & 76,25 & 50,0--95,0 \\\hline
\end{tabular}
\legend{Fonte: dados da pesquisa (2026).}
\end{table}

A Tabela~\ref{tab:ap-ueqs} traz os oito pares do UEQ-S em valores brutos, de 1 a 7, com o polo negativo à esquerda. Os quatro primeiros formam a dimensão pragmática e os quatro últimos, a hedônica. Depois da subtração de 4, as médias do painel foram 1,78 na dimensão pragmática, 1,88 na hedônica e 1,83 no valor geral.

\begin{table}[htbp]
\centering
\caption{Pares do UEQ-S em valores brutos (escala de 1 a 7).}
\label{tab:ap-ueqs}
\footnotesize
\renewcommand{\arraystretch}{1.2}
\begin{tabular}{|p{8.4cm}|c|c|c|c|}
\hline
\textbf{Par} & \textbf{n} & \textbf{Média} & \textbf{Mediana} & \textbf{Mín.--máx.} \\\hline
A plataforma atrapalhou ou apoiou o que você queria fazer? & 8 & 6,00 & 6,5 & 2--7 \\\hline
A plataforma é complicada ou fácil de usar? & 8 & 5,62 & 6 & 3--7 \\\hline
A plataforma é ineficiente ou eficiente? & 8 & 5,88 & 6,5 & 2--7 \\\hline
A plataforma é confusa ou clara? & 8 & 5,62 & 6 & 3--7 \\\hline
A plataforma é entediante ou empolgante? & 8 & 6,12 & 6,5 & 4--7 \\\hline
A plataforma é desinteressante ou interessante? & 8 & 6,25 & 6,5 & 5--7 \\\hline
A plataforma é convencional ou inventiva? & 8 & 5,38 & 6 & 1--7 \\\hline
A plataforma é comum ou inovadora? & 8 & 5,75 & 6 & 3--7 \\\hline
\end{tabular}
\legend{Fonte: dados da pesquisa (2026).}
\end{table}

A Tabela~\ref{tab:ap-qualidade} traz o bloco de qualidade técnica, e a Tabela~\ref{tab:ap-multipla}, as perguntas de múltipla escolha que não estão no Capítulo~\ref{cap:resultados}.

\begin{table}[htbp]
\centering
\caption{Bloco de qualidade técnica (escala de 1 a 5; ``não sei'' contado à parte).}
\label{tab:ap-qualidade}
\footnotesize
\renewcommand{\arraystretch}{1.2}
\begin{tabular}{|p{7.6cm}|c|c|c|c|c|}
\hline
\textbf{Afirmação} & \textbf{n} & \textbf{Média} & \textbf{Mediana} & \textbf{Mín.--máx.} & \textbf{Não sei} \\\hline
As mensagens de erro são claras e indicam o que precisa ser corrigido. & 7 & 4,71 & 5 & 4--5 & 1 \\\hline
A ferramenta respondeu de forma rápida e estável durante o uso. & 8 & 4,88 & 5 & 4--5 & 0 \\\hline
O assistente de criação de linguagem é compreensível para quem não tem formação em compiladores. & 8 & 4,38 & 4,5 & 3--5 & 0 \\\hline
A validação prévia impede adequadamente configurações de linguagem inválidas. & 6 & 4,67 & 5 & 4--5 & 2 \\\hline
O código intermediário exibido corresponde ao que o programa realmente executa. & 6 & 4,83 & 5 & 4--5 & 2 \\\hline
\end{tabular}
\legend{Fonte: dados da pesquisa (2026).}
\end{table}

\begin{table}[htbp]
\centering
\caption{Perguntas de múltipla escolha do formulário dos especialistas.}
\label{tab:ap-multipla}
\footnotesize
\renewcommand{\arraystretch}{1.2}
\begin{tabular}{|p{10.5cm}|c|}
\hline
\textbf{Resposta} & \textbf{Respondentes} \\\hline
\multicolumn{2}{|l|}{\textbf{Construções que fizeram falta (mais de uma marcação; sete respondentes)}} \\\hline
Não senti falta de nada relevante para o ensino introdutório & 5 \\\hline
Registros ou estruturas compostas & 2 \\\hline
Modularização em vários arquivos & 1 \\\hline
Orientação a objetos & 1 \\\hline
Manipulação de cadeias mais rica & 1 \\\hline
\multicolumn{2}{|l|}{\textbf{Momento da disciplina em que a plataforma seria mais útil (só quem leciona)}} \\\hline
Em paralelo à linguagem oficial, como apoio & 2 \\\hline
Não vejo espaço para ela na disciplina & 1 \\\hline
Antes da linguagem oficial da disciplina, como preparação & 1 \\\hline
\multicolumn{2}{|l|}{\textbf{Autorização de citação das respostas abertas}} \\\hline
Sim, autorizo a citação sem identificação & 7 \\\hline
Não autorizo & 1 \\\hline
\end{tabular}
\legend{Fonte: dados da pesquisa (2026).}
\end{table}

\section{Formulário da Sessão com Estudantes}

O formulário recebeu dezoito respostas em 3 de novembro de 2026, e todas trazem a opção ``Li e quero participar'' no termo de assentimento. A Tabela~\ref{tab:ap-estudantes-perfil} traz o perfil da turma, a Tabela~\ref{tab:ap-estudantes-dist} mostra quantos alunos marcaram cada ponto da escala em cada afirmação e a Tabela~\ref{tab:resultado-versoes} traz a comparação entre as duas versões do programa. O campo de comentário livre recebeu duas respostas, que não são transcritas porque o formulário não pede autorização de citação.

\begin{table}[htbp]
\centering
\caption{Perfil dos estudantes (dezoito respondentes).}
\label{tab:ap-estudantes-perfil}
\footnotesize
\renewcommand{\arraystretch}{1.2}
\begin{tabular}{|p{10.5cm}|c|}
\hline
\textbf{Resposta} & \textbf{Estudantes} \\\hline
\multicolumn{2}{|l|}{\textbf{Idade}} \\\hline
Até 15 anos & 9 \\\hline
16 ou 17 anos & 9 \\\hline
\multicolumn{2}{|l|}{\textbf{Curso técnico}} \\\hline
Informática & 18 \\\hline
\multicolumn{2}{|l|}{\textbf{Ano}} \\\hline
1º ano & 18 \\\hline
\multicolumn{2}{|l|}{\textbf{Experiência com programação}} \\\hline
Já programei um pouco (em aulas, cursos ou tutoriais) & 9 \\\hline
Programo com frequência & 8 \\\hline
Nunca programei & 1 \\\hline
\end{tabular}
\legend{Fonte: dados da pesquisa (2026).}
\end{table}

\begin{table}[htbp]
\centering
\caption{Número de estudantes em cada ponto da escala, de 1 (discordo totalmente) a 5 (concordo totalmente).}
\label{tab:ap-estudantes-dist}
\footnotesize
\renewcommand{\arraystretch}{1.2}
\begin{tabular}{|p{7.4cm}|c|c|c|c|c|c|}
\hline
\textbf{Afirmação} & \textbf{1} & \textbf{2} & \textbf{3} & \textbf{4} & \textbf{5} & \textbf{Média} \\\hline
Entendi para que serve a plataforma. & 1 & 0 & 0 & 8 & 9 & 4,33 \\\hline
Um programa escrito com palavras em português é mais fácil de entender do que um escrito com palavras em inglês. & 0 & 1 & 6 & 3 & 8 & 4,00 \\\hline
Poder escolher as palavras da linguagem ajudaria quem está começando a programar. & 0 & 1 & 3 & 6 & 8 & 4,17 \\\hline
Ver como a plataforma separa o programa em partes (os \textit{tokens}) ajuda a entender como o computador lê o código. & 0 & 2 & 1 & 7 & 8 & 4,17 \\\hline
Eu gostaria de usar essa plataforma nas aulas de programação. & 0 & 0 & 3 & 9 & 6 & 4,17 \\\hline
A demonstração aumentou meu interesse em aprender programação. & 1 & 1 & 0 & 8 & 8 & 4,17 \\\hline
\end{tabular}
\legend{Fonte: dados da pesquisa (2026).}
\end{table}
